% q0054.tex — Bloppar Blues
\question{
  id        = {0054},
  title     = {Bloppar Blues},
  source    = {OBECON},
  year      = {2026},
  round     = {Seletiva IEO -- Phase C},
  language  = {English},
  topic     = {Macro},
  subtopic  = {Exchange Rate / Balance Sheet Effects},
  type      = {objective},
  statement = {%
    A Xieconomian firm exports all of its output (earning Bloppars), has all
    its operating costs denominated in Rons, holds Ron-denominated reserves,
    but carries Bloppar-denominated debt. The Ron unexpectedly appreciates
    against the Bloppar. Standard analysis suggests appreciation hurts
    exporters, yet the firm's equity (measured in Rons) increases. Why?
  },
  choices   = {%
    \item Appreciation increases foreign demand for exports, making them
          cheaper in Bloppars.
    \item Appreciation raises the Ron value of the firm's Ron reserves,
          offsetting export revenue losses.
    \item Appreciation improves net worth by reducing the Ron value of
          foreign-currency liabilities.
    \item Appreciation lowers imported input costs, improving profit margins.
  },
  answer    = {C},
  solution  = {%
    Two effects arise from the Ron appreciation. First, operating profit
    suffers: Bloppar export revenues convert to fewer Rons. Second, a
    balance-sheet effect: Bloppar-denominated liabilities become cheaper when
    measured in Rons. Because the firm carries significant Bloppar debt, the
    valuation gain from the reduced debt burden (in Ron terms) dominates the
    revenue loss, increasing shareholder equity. (C) is correct.
  },
}
